\documentclass[10pt]{beamer}

%%%%%%%%%%%%%%%%%%%%%%%%%%%%%%%%%%%%%%%%%%%%%%%%%%%%%%%%%%%%%%%%%%
%                        PACKAGES                                %
%%%%%%%%%%%%%%%%%%%%%%%%%%%%%%%%%%%%%%%%%%%%%%%%%%%%%%%%%%%%%%%%%%

\usepackage[utf8]{inputenc}
\usepackage[T1]{fontenc}
\usepackage[french]{babel}
\usepackage{amsmath,amssymb,amsthm,mathtools}
\usepackage{booktabs}
\usepackage{array}
\usepackage{tikz}
\usetikzlibrary{arrows.meta, positioning, calc}
\usepackage{pgfplots}
\pgfplotsset{compat=1.16}

%%%%%%%%%%%%%%%%%%%%%%%%%%%%%%%%%%%%%%%%%%%%%%%%%%%%%%%%%%%%%%%%%%
%                     THÈME BEAMER                               %
%%%%%%%%%%%%%%%%%%%%%%%%%%%%%%%%%%%%%%%%%%%%%%%%%%%%%%%%%%%%%%%%%%

\definecolor{mainblue}{RGB}{20, 60, 120}
\definecolor{lightblue}{RGB}{220, 232, 248}
\definecolor{midblue}{RGB}{100, 140, 200}
\definecolor{darktext}{RGB}{30, 30, 30}

\usetheme{Madrid}
\usecolortheme[named=midblue]{structure}
\useoutertheme{infolines}
\setbeamertemplate{navigation symbols}{}
\setbeamertemplate{itemize item}{\small$\bullet$}
\setbeamertemplate{itemize subitem}{\small$\circ$}

\setbeamercolor{block title}{bg=mainblue, fg=white}
\setbeamercolor{block body}{bg=lightblue, fg=darktext}
\setbeamercolor{block title example}{bg=mainblue!70!black, fg=white}
\setbeamercolor{block body example}{bg=lightblue!60, fg=darktext}

\setbeamertemplate{footline}{
  \leavevmode\hbox{%
    \begin{beamercolorbox}[wd=.33\paperwidth,ht=2.25ex,dp=1ex,center]{author in head/foot}%
      \usebeamerfont{author in head/foot}\insertshortauthor
    \end{beamercolorbox}%
    \begin{beamercolorbox}[wd=.34\paperwidth,ht=2.25ex,dp=1ex,center]{title in head/foot}%
      \usebeamerfont{title in head/foot}\insertshorttitle
    \end{beamercolorbox}%
    \begin{beamercolorbox}[wd=.33\paperwidth,ht=2.25ex,dp=1ex,right]{date in head/foot}%
      \usebeamerfont{date in head/foot}
      \insertframenumber{} / \inserttotalframenumber\hspace*{1.5ex}
    \end{beamercolorbox}}%
  \vskip0pt%
}

\newcommand{\framewithtitle}[1]{%
\begin{frame}[c]{}
  \begin{center}
  \Huge
    \begin{tcolorbox}
      \emph{
        \espace
        \begin{center}
          \textcolor{black}{#1}
        \end{center}
        \espace
      }
    \end{tcolorbox}
  \end{center}
\end{frame}
}

%%%%%%%%%%%%%%%%%%%%%%%%%%%%%%%%%%%%%%%%%%%%%%%%%%%%%%%%%%%%%%%%%%
%                    TITRE / AUTEURS                             %
%%%%%%%%%%%%%%%%%%%%%%%%%%%%%%%%%%%%%%%%%%%%%%%%%%%%%%%%%%%%%%%%%%

\title[Équation de Nagumo]{Analyse et simulation numérique\\de l'équation de Nagumo}
\subtitle{Étude théorique, schémas numériques et stabilité}
\titlegraphic{\includegraphics[height=2.0cm]{img/logo1.jpg}}
\author[GLODJO \& HOMEVO \& AGRO]{M.\,J. GLODJO \quad K. HOMEVO \quad  T. Rachidi \\[4pt]
  \small Sous la supervision de \textbf{Dr Jamal ADETOLA}}
\institute[ENSGMM -- UNSTIM]{TP de Problèmes instationnaires -- ENSGMM}
\date{}

%%%%%%%%%%%%%%%%%%%%%%%%%%%%%%%%%%%%%%%%%%%%%%%%%%%%%%%%%%%%%%%%%%
\begin{document}
%%%%%%%%%%%%%%%%%%%%%%%%%%%%%%%%%%%%%%%%%%%%%%%%%%%%%%%%%%%%%%%%%%

%--- Diapo 1 : Page de titre
\begin{frame}
  \titlepage
\end{frame}

%--- Diapo 2 : Plan
\begin{frame}{Plan}
  \tableofcontents[hideallsubsections]
\end{frame}

\begin{frame}{Introduction}
\begin{itemize}
    \item \textbf{Contexte} : circuit électronique actif proposé par Nagumo, Arimoto et Yoshizawa (1962), prolongeant les travaux de Hodgkin \& Huxley.
    \item \textbf{Réduction bistable} du système de Hodgkin-Huxley $\rightarrow$ équation de réaction-diffusion non linéaire (équation de Nagumo).
    \item \textbf{Non-linéarité cubique} $\rightarrow$ structure bistable : deux états stables séparés par un seuil instable.
    \item \textbf{Propriété clé} : génère des fronts voyageurs reliant ces deux états – modèle paradigmatique pour l'étude analytique et numérique.
    \item \textbf{Objectifs du rapport} :
    \begin{itemize}
        \item Analyse théorique rigoureuse des propriétés qualitatives.
        \item Construction et analyse de quatre schémas numériques.
        \item Étude de la stabilité des schémas et établissement des conditions de convergence.
        \item Validation par simulation numérique et comparaison des performances.
    \end{itemize}
\end{itemize}
\end{frame}

\section{Analyse théorique}
\framewithtitle{Analyse théorique}

\begin{frame}{Modèle générique de réaction-diffusion}

\begin{equation}
  \frac{\partial u}{\partial t} = D\,\Delta u + f(u)
\end{equation}

\bigskip

\textbf{Forme de Nagumo} — cinétique cubique bistable :
\begin{equation}
  \frac{\partial u}{\partial t} = D\,\Delta u + u(1-u)(u-a), \qquad a \in (0,1)
\end{equation}

\bigskip

\begin{tabular}{@{}lll@{}}
  \toprule
  Symbole & Nom & Rôle \\
  \midrule
  $u(\mathbf{x,y},t)$ & Potentiel de membrane   & Variable d'état adimensionnée \\[4pt]
  $D > 0$           & Coefficient de diffusion & Propagation passive par couplage électrotonique \\[4pt]
  $f(u)$            & Terme de réaction        & Mécanismes ioniques d'excitabilité \\[4pt]
  $a \in (0,1)$     & Seuil d'excitabilité     & Contrôle le basculement entre les deux états stables \\
  \bottomrule
\end{tabular}

\end{frame}

\begin{frame}{Formulation du problème}

On cherche $u : \Omega \times (0,+\infty) \to \mathbb{R}$ solution de :

\begin{equation}
  \begin{cases}
    \displaystyle\frac{\partial u}{\partial t} = D\,\Delta u + u(1-u)(u-a),
    & (x,y) \in \Omega,\; t > 0, \\[10pt]
    u(x,y,0) = u_0(x,y),
    & (x,y) \in \Omega, \\[6pt]
    \text{C.L. sur } \partial\Omega,
    & t > 0.
  \end{cases}
\end{equation}

\bigskip

\textbf{Conditions aux limites usuelles}

\medskip

\begin{tabular}{@{}lll@{}}
  \toprule
  Type & Condition & Interprétation physique \\
  \midrule
  Dirichlet & $u = 0 \text{ sur } \partial\Omega$ & Bord maintenu au potentiel de repos \\[4pt]
  Neumann   & $\dfrac{\partial u}{\partial \mathbf{n}} = 0 \text{ sur } \partial\Omega$
            & Isolation électrique — flux nul au bord \\
  \bottomrule
\end{tabular}

\end{frame}

\begin{frame}{Équilibres et stabilité}

Les états homogènes stationnaires vérifient $f(u^*) = 0$ :

\begin{equation}
  u^*(1-u^*)(u^*-a) = 0
  \qquad \Longrightarrow \qquad
  u^*_1 = 0, \quad u^*_2 = a, \quad u^*_3 = 1.
\end{equation}

\bigskip

\textbf{Analyse de stabilité linéaire} — signe de $f'(u^*)$ :

\begin{equation}
  f'(u) = -3u^2 + 2(1+a)u - a
\end{equation}

\medskip

\begin{tabular}{@{}llll@{}}
  \toprule
  Équilibre & $f'(u^*)$ & Stabilité & Rôle \\
  \midrule
  $u^* = 0$ & $-a < 0$        & asymptotiquement stable & état de repos \\[4pt]
  $u^* = a$ & $a(1-a) > 0$    & instable                & seuil d'excitation \\[4pt]
  $u^* = 1$ & $-(1-a) < 0$    & asymptotiquement stable & état excité \\
  \bottomrule
\end{tabular}

\end{frame}


\begin{frame}{Fronts voyageurs en 2D}

Ansatz : $u(x,y,t) = \phi(\xi)$, \quad $\xi = x\cos\theta + y\sin\theta - ct$

\medskip

\textbf{ODE du front} — identique au cas 1D :
\begin{equation}
  D\,\phi'' + c\,\phi' + \phi(1-\phi)(\phi-a) = 0,
  \qquad \phi(-\infty)=1,\quad \phi(+\infty)=0.
\end{equation}

\textbf{Solution exacte et vitesse critique :}
\begin{equation}
  \phi^*(\xi) = \frac{1}{1+e^{\,\xi/\sqrt{2D}}}, \qquad
  \boxed{c^* = \sqrt{2D}\!\left(\tfrac{1}{2}-a\right).}
\end{equation}

\textbf{Propriétés clés :}

\begin{itemize}
  \item $c^*$ est indépendante de la direction de propagation $\theta$
  \item $a < \tfrac{1}{2}$ : front avançant ; \quad
        $a = \tfrac{1}{2}$ : front stationnaire ; \quad
        $a > \tfrac{1}{2}$ : front reculant
  \item largeur caractéristique du profil : $\ell = \sqrt{2D}$
  \item stabilité orbitale : un front perturbé converge vers $\phi^*$
\end{itemize}

\end{frame}

\begin{frame}{Principe du maximum et existence globale}

Pour $u_0 \in L^\infty(\Omega)$ et $\Omega$ borné régulier :

\bigskip

\textbf{Principe du maximum :}

\begin{itemize}
  \item $u_0(x,y) \in [0,1] \;\Rightarrow\; u(x,y,t) \in [0,1]$ pour tout $t \geq 0$
  \item au point de maximum $\Delta u \leq 0$, donc si $M > 1$ alors $f(M) < 0$
        et $\partial_t M < 0$ : le maximum ne peut dépasser~1
\end{itemize}

\bigskip

\textbf{Conséquences :}

\begin{itemize}
  \item borne a priori $\Rightarrow$ existence globale sur $[0, +\infty)$
  \item la solution reste dans la région physique $[0,1]$
  \item les schémas explicites avec $\Delta t$ trop grand violent ce principe
        et produisent des oscillations parasites
\end{itemize}

\end{frame}

\section{Analyse numérique}
\framewithtitle{Analyse numérique}

\begin{frame}{Discrétisation spatiale — schéma à cinq points}

Grille uniforme : $x_i = i\,\Delta x$, $y_j = j\,\Delta y$, \quad
$u_{i,j}^n \approx u(x_i, y_j, t^n)$.

\bigskip

\begin{center}
\begin{tikzpicture}[scale=1.2]
  % Grille légère
  \draw[gray!30] (-2,-2) grid (2,2);

  % Liaisons du stencil
  \draw[thick] (0,0) -- (1,0);
  \draw[thick] (0,0) -- (-1,0);
  \draw[thick] (0,0) -- (0,1);
  \draw[thick] (0,0) -- (0,-1);

  % Points voisins
  \filldraw[black] ( 1, 0) circle (3pt) node[right=4pt] {$u_{i+1,j}$};
  \filldraw[black] (-1, 0) circle (3pt) node[left=4pt]  {$u_{i-1,j}$};
  \filldraw[black] ( 0, 1) circle (3pt) node[above=4pt] {$u_{i,j+1}$};
  \filldraw[black] ( 0,-1) circle (3pt) node[below=4pt] {$u_{i,j-1}$};

  % Point central
  \filldraw[red] (0,0) circle (3.5pt) node[above right=4pt] {$u_{i,j}$};

  % Axes
  \draw[->] (-2.2,0) -- (2.4,0) node[right] {$x$};
  \draw[->] (0,-2.2) -- (0,2.4) node[above] {$y$};

  % Labels dx, dy
  \draw[<->] (0,-1.6) -- (1,-1.6) node[midway, below] {$\Delta x$};
  \draw[<->] (1.6,0)  -- (1.6,1) node[midway, right] {$\Delta y$};
\end{tikzpicture}
\end{center}

\end{frame}

\begin{frame}{Discrétisation spatiale — détails}

\textbf{Approximation du Laplacien 2D :}
\begin{equation}
  \Delta_h u_{i,j}^n =
  \frac{u_{i+1,j}^n - 2u_{i,j}^n + u_{i-1,j}^n}{(\Delta x)^2}
  + \frac{u_{i,j+1}^n - 2u_{i,j}^n + u_{i,j-1}^n}{(\Delta y)^2}
\end{equation}

\bigskip

\textbf{Propriétés de $A_{2D}$ :}

\begin{itemize}
  \item structure de Kronecker :
        $A_{2D} = I_{N_y} \otimes A_x + A_y \otimes I_{N_x}$
  \item matrice creuse, symétrique, semi-définie négative
  \item valeur propre extrême : $|\lambda_{\max}| = 8D/h^2$
        — deux fois plus grande qu'en 1D
  \item erreur de troncature : $O\!\left((\Delta x)^2 + (\Delta y)^2\right)$
\end{itemize}

\bigskip

\textbf{Conditions de Neumann} — points fantômes :
\begin{equation}
  u_{-1,j}^n = u_{1,j}^n, \quad
  u_{N_x+1,j}^n = u_{N_x-1,j}^n, \quad
  u_{i,-1}^n = u_{i,1}^n, \quad
  u_{i,N_y+1}^n = u_{i,N_y-1}^n.
\end{equation}

\end{frame}

\begin{frame}{Runge-Kutta 4 — Description de l'algorithme}

Méthode explicite à un pas, d'ordre 4 en temps. On note
$G(\mathbf{U}) = A_{2D}\mathbf{U} + F(\mathbf{U})$ le flux total.

\bigskip

\textbf{Les quatre étapes :}

\begin{itemize}
  \item $K_1 = G(\mathbf{U}^n)$
  \item $K_2 = G\!\left(\mathbf{U}^n + \tfrac{\Delta t}{2}K_1\right)$
  \item $K_3 = G\!\left(\mathbf{U}^n + \tfrac{\Delta t}{2}K_2\right)$
  \item $K_4 = G\!\left(\mathbf{U}^n + \Delta t\, K_3\right)$
\end{itemize}

\bigskip

\textbf{Mise \`a jour :}
\begin{equation}
  \boxed{\mathbf{U}^{n+1} = \mathbf{U}^n
    + \frac{\Delta t}{6}\bigl(K_1 + 2K_2 + 2K_3 + K_4\bigr).}
\end{equation}

\end{frame}

\begin{frame}{Runge-Kutta 4 — Schéma numérique}

\textbf{Mise à jour :}
\begin{equation}
  \boxed{\mathbf{U}^{n+1} = \mathbf{U}^n
    + \frac{\Delta t}{6}\bigl(K_1 + 2K_2 + 2K_3 + K_4\bigr)}
\end{equation}

\bigskip

\textbf{Implémentation pratique du produit $A_{2D}\mathbf{U}$ :}
\begin{equation}
  (A_{2D}\mathbf{U})_{i,j} =
  D\,\frac{U_{i-1,j} - 2U_{i,j} + U_{i+1,j}}{(\Delta x)^2}
  + D\,\frac{U_{i,j-1} - 2U_{i,j} + U_{i,j+1}}{(\Delta y)^2}
\end{equation}

\medskip

\begin{itemize}
  \item exploitation de la structure à cinq points — pas de stockage de $A_{2D}$
  \item $\mathbf{U}^n$ stocké comme tableau 2D de taille $(N_x{+}1)\times(N_y{+}1)$
  \item 4 évaluations de $G$ par pas de temps, coût $O(N_x N_y)$ chacune
\end{itemize}

\end{frame}

\begin{frame}{Runge-Kutta 4 — Stabilité et convergence}

\textbf{Facteur d'amplification} sur $u' = \lambda u$, $z = \Delta t\,\xi_{k_1,k_2}$ :
\begin{equation}
  \zeta(z) = 1 + z + \frac{z^2}{2} + \frac{z^3}{6} + \frac{z^4}{24}
\end{equation}

\bigskip

\textbf{Condition de stabilité CFL en 2D} — $|\zeta(z)| \leq 1$ pour $z \geq -2{,}79$ :
\begin{equation}
  \Delta t \cdot \frac{8D}{h^2} \leq 2{,}79
  \qquad \Longrightarrow \qquad
  \boxed{\mu = \frac{D\,\Delta t}{h^2} \leq 0{,}349}
\end{equation}

\bigskip

\begin{itemize}
  \item ordre 4 en temps, ordre 2 en espace — erreur globale $O((\Delta t)^4 + h^2)$
  \item condition CFL \textbf{deux fois plus restrictive} qu'en 1D ($\mu \leq 0{,}697$)
  \item stabilité conditionnelle : $\Delta t$ contraint par le maillage
  \item meilleure précision temporelle parmi les schémas explicites
\end{itemize}

\end{frame}

\begin{frame}{Euler Implicite — Description de l'algorithme}

Approche \textbf{semi-implicite} : diffusion traitée implicitement,
réaction traitée explicitement.

\bigskip

\textbf{Principe :}
\begin{itemize}
  \item discrétisation temporelle par différence progressive d'ordre 1
  \item terme diffusif évalué en $t^{n+1}$ — stabilité inconditionnelle
  \item terme réactif évalué en $t^n$ — linéarité du système à résoudre
\end{itemize}

\bigskip

\textbf{Algorithme par pas de temps :}
\begin{itemize}
  \item calculer le second membre $\mathbf{b}^n = \mathbf{U}^n + \Delta t\,F(\mathbf{U}^n)$
  \item résoudre le système linéaire $(I - \Delta t\,A_{2D})\,\mathbf{U}^{n+1} = \mathbf{b}^n$
  \item pour les grandes grilles : utiliser la méthode \textbf{ADI de Peaceman-Rachford}
        qui décompose le système 2D en systèmes 1D tridiagonaux
\end{itemize}

\end{frame}

\begin{frame}{Euler Implicite — Schéma numérique}

\textbf{Système linéaire global :}
\begin{equation}
  \boxed{(I - \Delta t\,A_{2D})\,\mathbf{U}^{n+1} = \mathbf{U}^n + \Delta t\,F(\mathbf{U}^n)}
\end{equation}

\bigskip

\textbf{Méthode ADI de Peaceman-Rachford :}

\medskip

Demi-pas 1 — implicite en $x$, explicite en $y$ :
\begin{equation}
  \left(I - \tfrac{\Delta t}{2}D\,A_x\right)\mathbf{U}^{n+1/2}
  = \left(I + \tfrac{\Delta t}{2}D\,A_y\right)\mathbf{U}^n + \Delta t\,F(\mathbf{U}^n)
\end{equation}

Demi-pas 2 — implicite en $y$, explicite en $x$ :
\begin{equation}
  \left(I - \tfrac{\Delta t}{2}D\,A_y\right)\mathbf{U}^{n+1}
  = \left(I + \tfrac{\Delta t}{2}D\,A_x\right)\mathbf{U}^{n+1/2}
\end{equation}

\medskip

\begin{itemize}
  \item chaque demi-pas réduit à $N_y{+}1$ (resp. $N_x{+}1$) systèmes tridiagonaux 1D
  \item coût total par pas : $O(N_x N_y)$ au lieu de $O((N_x N_y)^{1{,}5})$
\end{itemize}

\end{frame}

\begin{frame}{Euler Implicite — Stabilité et convergence}

\textbf{Facteur d'amplification} — mode de Fourier $z = \Delta t\,\xi_{k_1,k_2} \leq 0$ :
\begin{equation}
  \rho = \frac{1}{1 - z} = \frac{1}{1 + |\Delta t\,\xi_{k_1,k_2}|} \in (0,1]
  \qquad \forall\, \mu > 0.
\end{equation}

\bigskip

\begin{itemize}
  \item \textbf{inconditionnellement stable} — aucune contrainte sur $\Delta t$
  \item ordre 1 en temps, ordre 2 en espace — erreur globale $O(\Delta t + h^2)$
  \item matrice $(I - \Delta t\,A_{2D})$ symétrique définie positive — système
        toujours inversible
  \item avec réaction explicite ($\lambda = f'(u^*)$) :
  \begin{equation}
    \rho = \frac{1 + \Delta t\,\lambda}{1 - \Delta t\,\xi_{k_1,k_2}}
  \end{equation}
  stabilité maintenue aux équilibres stables ($\lambda < 0$)
\end{itemize}

\end{frame}

\begin{frame}{Adams-Bashforth 2 — Description de l'algorithme}

Schéma \textbf{multipass explicite} d'ordre 2. On approche l'intégrale exacte
$\mathbf{U}^{n+1} - \mathbf{U}^n = \int_{t^n}^{t^{n+1}} G(\mathbf{U})\,dt$
par interpolation de Lagrange sur deux niveaux passés.

\bigskip

\textbf{Principe :}
\begin{itemize}
  \item interpoler $G$ sur $(t^{n-1}, G^{n-1})$ et $(t^n, G^n)$
        par un polynôme de degré 1
  \item intégrer exactement ce polynôme sur $[t^n, t^{n+1}]$
        $\Rightarrow$ coefficients $\tfrac{3}{2}$ et $-\tfrac{1}{2}$
  \item réutiliser $G^n$ au pas suivant comme $G^{n-1}$ — un seul calcul de $G$ par pas
\end{itemize}

\bigskip

\textbf{Initialisation :}
\begin{itemize}
  \item $\mathbf{U}^0$ : condition initiale
  \item $\mathbf{U}^1$ : un pas d'Euler explicite
        $\mathbf{U}^1 = \mathbf{U}^0 + \Delta t\,G(\mathbf{U}^0)$
\end{itemize}

\end{frame}

\begin{frame}{Adams-Bashforth 2 — Schéma numérique}

\textbf{Schéma AB2 en chaque point $(i,j)$ :}
\begin{equation}
  \boxed{u_{i,j}^{n+1} = u_{i,j}^n
    + \Delta t\!\left(\frac{3}{2}\,G_{i,j}^n - \frac{1}{2}\,G_{i,j}^{n-1}\right)}
\end{equation}

avec le flux total :
\begin{equation}
  G_{i,j}^n = D\!\left(
    \frac{u_{i+1,j}^n - 2u_{i,j}^n + u_{i-1,j}^n}{(\Delta x)^2}
  + \frac{u_{i,j+1}^n - 2u_{i,j}^n + u_{i,j-1}^n}{(\Delta y)^2}
  \right) + f(u_{i,j}^n)
\end{equation}

\bigskip

\textbf{Organisation mémoire :}
\begin{itemize}
  \item stocker $G^{n-1}$ et $G^n$ entre deux pas consécutifs
  \item à chaque pas : calculer $\mathbf{U}^{n+1}$, puis mettre à jour
        $G^{n-1} \leftarrow G^n$, $G^n \leftarrow G(\mathbf{U}^{n+1})$
\end{itemize}

\end{frame}

\begin{frame}{Adams-Bashforth 2 — Stabilité et convergence}

\textbf{Polynôme caractéristique} — mode $z = \Delta t\,\xi_{k_1,k_2} \leq 0$ :
\begin{equation}
  \zeta^2 - \left(1 + \frac{3}{2}z\right)\zeta + \frac{z}{2} = 0
\end{equation}

\bigskip

\textbf{Condition de stabilité CFL en 2D} — région $z \in [-1, 0]$ :
\begin{equation}
  \Delta t \cdot \frac{8D}{h^2} \leq 1
  \qquad \Longrightarrow \qquad
  \boxed{\mu = \frac{D\,\Delta t}{h^2} \leq \frac{1}{8}}
\end{equation}

\bigskip

\begin{itemize}
  \item ordre 2 en temps, ordre 2 en espace — erreur globale $O((\Delta t)^2 + h^2)$
  \item CFL \textbf{deux fois plus restrictive} qu'en 1D ($\mu \leq \tfrac{1}{4}$)
  \item CFL \textbf{quatre fois plus restrictive} qu'Euler explicite en 1D
  \item un seul calcul de $G$ par pas — coût $O(N_x N_y)$
\end{itemize}

\end{frame}

\begin{frame}{Adams-Moulton 2 — Description de l'algorithme}

Schéma \textbf{semi-implicite} d'ordre 2, basé sur la règle des trapèzes.
Même stratégie qu'Euler implicite : diffusion implicite, réaction explicite.

\bigskip

\textbf{Principe :}
\begin{itemize}
  \item approximer $\int_{t^n}^{t^{n+1}} G\,dt$ par la règle des trapèzes
  \item traiter $A_{2D}\mathbf{U}$ implicitement aux deux instants $t^n$ et $t^{n+1}$
  \item approximer la réaction explicitement par $2F(\mathbf{U}^n)$
        pour éviter toute résolution non linéaire
\end{itemize}

\bigskip

\textbf{Résolution par ADI de Peaceman-Rachford :}
\begin{itemize}
  \item demi-pas 1 : implicite en $x$, explicite en $y$
  \item demi-pas 2 : implicite en $y$, explicite en $x$
  \item chaque demi-pas réduit à des systèmes tridiagonaux 1D
\end{itemize}

\end{frame}

\begin{frame}{Adams-Moulton 2 — Schéma numérique}

\textbf{Système linéaire global :}
\begin{equation}
  \boxed{\left(I - \frac{\Delta t}{2}A_{2D}\right)\mathbf{U}^{n+1}
  = \mathbf{U}^n + \frac{\Delta t}{2}\!\left(A_{2D}\mathbf{U}^n + 2F(\mathbf{U}^n)\right)}
\end{equation}

\bigskip

\textbf{Décomposition ADI :}

\medskip

Demi-pas 1 — implicite en $x$ :
\begin{equation}
  \left(I - \tfrac{\Delta t}{2}D\,A_x\right)\mathbf{U}^{n+1/2}
  = \left(I + \tfrac{\Delta t}{2}D\,A_y\right)\mathbf{U}^n + \Delta t\,F(\mathbf{U}^n)
\end{equation}

Demi-pas 2 — implicite en $y$ :
\begin{equation}
  \left(I - \tfrac{\Delta t}{2}D\,A_y\right)\mathbf{U}^{n+1}
  = \left(I + \tfrac{\Delta t}{2}D\,A_x\right)\mathbf{U}^{n+1/2}
\end{equation}

\end{frame}

\begin{frame}{Adams-Moulton 2 — Stabilité et convergence}

\textbf{Facteur d'amplification} — $z = \Delta t\,\xi_{k_1,k_2} \leq 0$ :
\begin{equation}
  \zeta = \frac{1 + z/2}{1 - z/2}
\end{equation}

\textbf{Démonstration d'A-stabilité :}
\begin{equation}
  (1 - z/2)^2 - (1 + z/2)^2 = -2z \geq 0 \quad \text{pour } z \leq 0
  \quad \Longrightarrow \quad |\zeta| \leq 1 \quad \forall\, z \leq 0.
\end{equation}

\bigskip

\begin{itemize}
  \item \textbf{inconditionnellement stable} (A-stable) — aucune contrainte sur $\Delta t$
  \item ordre 2 en temps, ordre 2 en espace — erreur globale $O((\Delta t)^2 + h^2)$
  \item $|\zeta| < 1$ strictement pour $z < 0$ — propriété plus forte que la stabilité
        sur l'axe réel négatif
  \item avec l'ADI : coût $O(N_x N_y)$ par pas — \textbf{meilleur compromis ordre/coût}
        pour les grandes grilles 2D
\end{itemize}

\end{frame}

\begin{frame}{Comparaison des quatre schémas en 2D}

\begin{tabular}{@{}lllll@{}}
  \toprule
  Méthode & Type & Ordre & CFL 2D ($\mu$) & Coût/pas \\
  \midrule
  RK4         & explicite & 4 & $\leq 0{,}349$      & $4 \times O(N^2)$ \\[4pt]
  Euler impl. & implicite & 1 & aucune              & $O(N^2)$ (ADI) \\[4pt]
  AB2         & explicite & 2 & $\leq \frac{1}{8}$  & $1 \times O(N^2)$ \\[4pt]
  AM2 + ADI   & implicite & 2 & aucune              & $O(N^2)$ (ADI) \\
  \bottomrule
\end{tabular}

\bigskip

\textbf{Points clés :}

\begin{itemize}
  \item les conditions CFL explicites en 2D sont exactement \textbf{deux fois plus
        restrictives} qu'en 1D, en raison du facteur $|\lambda_{\max}^{2D}| = 2\,|\lambda_{\max}^{1D}|$
  \item AB2 est le schéma explicite le plus contraignant ($\mu \leq 1/8$) —
        quatre fois plus restrictif qu'Euler explicite en 1D
  \item Euler implicite et AM2 sont inconditionnellement stables ;
        AM2 est préférable car d'ordre 2
  \item \textbf{AM2 + ADI} offre le meilleur compromis :
        ordre 2, A-stable, coût $O(N^2)$ par pas
\end{itemize}

\end{frame}

\begin{frame}{Le criquet pèlerin -- Un bistable naturel}
  \framesubtitle{\textit{Schistocerca gregaria} dans les plaines sahéliennes}

\textbf{Pourquoi ce cas d'étude ?}\\[0.4em]
      Le criquet pèlerin est l'un des rares organismes
      qui exhibe naturellement une \textbf{bistabilité dure}~:
      il existe sous \textbf{deux formes biologiquement distinctes}.

      \medskip
      \begin{itemize}
          \item Phase solitaire \quad $u \approx 0$ \\
        Individus discrets, verts, non migrateurs.\\
        Impact agricole \textbf{négligeable}.
        \item Phase grégaire \quad $u \approx 1$ \\
         Essaims de milliards d'individus, jaunes et noirs,\\
        migrant sur \textbf{des milliers de km}.\\
        Destruction totale des cultures en quelques heures.
      \end{itemize}

      \vspace{0.3em}
      \begin{exampleblock}{Seuil de grégarisation \quad $u = a$}
        En dessous~: retour à $u=0$. Au dessus~: bascule vers $u=1$.\\
        $\Rightarrow$ exactement un \textbf{équilibre instable}.
      \end{exampleblock}
  
\end{frame}

\section{Simulation et résultats}
\framewithtitle{Simulation et résultats}



\begin{frame}{Paramètres de la simulation}
  \begin{itemize}
    \item Domaine spatial : $[0,L_x]\times[0,L_y]$ avec $L_x = L_y = 60.0$.
    \item Grille : $N_x = N_y = 80$, soit $81\times81$ points.
    \item Espacement spatial : $\Delta x = \Delta y = 60/80 = 0.75$.
    \item Paramètres physiques :
      \begin{itemize}
        \item Diffusion $D = 1.0$,
        \item Seuil $a = 0.25$,
        \item Temps final $T = 20.0$.
      \end{itemize}
    \item Conditions initiales principales étudiées :
      \begin{itemize}
        \item front plan vertical,
        \item onde circulaire,
        \item front plan perturbé.
      \end{itemize}
    \item Vitesse théorique du front :
      $c^* = \sqrt{2D}\,(\tfrac12 - a) \approx 0.3536$.
  \end{itemize}
\end{frame}

\begin{frame}{Schémas numériques comparés}
  \begin{itemize}
    \item RK4 : Runge-Kutta 4 explicit.
    \item EI : Euler implicite semi-implicite.
    \item ADI : Peaceman-Rachford semi-implicite avec résolutions 1D.
    \item AB2 : Adams-Bashforth 2 explicite multi-pas.
    \item AM2 : IMEX Crank-Nicolson semi-implicite.
  \end{itemize}
  \begin{itemize}
    \item Schémas explicites nécessitent un petit pas temporel $\Delta t$.
    \item Schémas semi-implicites permettent un pas plus grand pour stabiliser la diffusion.
  \end{itemize}
\end{frame}

\begin{frame}{Comparaison des schémas à $t \approx 20.0$}
  \centering
  \includegraphics[width=0.95\textwidth]{img/comparaison_schemas_t_20.00.png}
  \caption{Distribution de $u(x,y)$ pour chaque schéma au temps final.}
\end{frame}

\begin{frame}{Interprétation : comparaison des schémas}
  \begin{itemize}
    \item Tous les schémas capturent la propagation d'un front entre $u\approx1$ et $u\approx0$.
    \item Les schémas implicites montrent une meilleure stabilité globale pour les grandes étapes temporelles.
    \item Le schéma RK4 propose une solution détaillée et peut servir de référence de haute précision.
    \item Les différences observées proviennent du traitement de la diffusion et du terme non-linéaire.
  \end{itemize}
\end{frame}

\begin{frame}{Suivi du front et vitesse numérique}
  \centering
  \includegraphics[width=0.95\textwidth]{img/front_tracking.png}
  \captionof{Position du front et vitesse numérique comparée à $c^*$.}
\end{frame}

\begin{frame}{Interprétation : suivi du front}
  \begin{itemize}
    \item La position du front évolue linéairement avec le temps, en accord avec l'hypothèse de propagation de type travelling wave.
    \item La droite théorique $x = L_x/2 + c^* t$ sert de référence de vitesse.
    \item Les schémas explicites et semi-implicites restent proches de la vitesse théorique, ce qui valide la simulation.
    \item Les fluctuations numériques sont plus visibles sur la vitesse calculée par différences finies, en particulier pour AB2 et ADI.
  \end{itemize}
\end{frame}

\begin{frame}{Profil 3D final du champ $u$}
  \centering
  \includegraphics[width=0.75\textwidth]{img/surface_3d_RK4.png}
  \captionof{figure}{Surface 3D du profil final $u(x,y)$ pour le schéma RK4.}
\end{frame}

\begin{frame}{Interprétation : surface 3D}
  \begin{itemize}
    \item Le profil final montre un front net entre les zones à $u\approx1$ et $u\approx0$.
    \item L'utilisation d'une grille régulière en 2D met en évidence la symétrie de la propagation.
    \item L'aspect 3D facilite l'observation du lissage dû à la diffusion et du maintien de la forme du front.
  \end{itemize}
\end{frame}

\end{document}
